% ===================================================================
%  TABELL Nr 5 — Huset och dess byggande.
%  Traduction 2026 d'apres texto/io, controlee sur texto/fr et
%  texto/es. Consigne : voir texto/sv/10-tabell-01.tex.
%
%  Le tableau 5 est un DIALOGUE, et les deux volumes composent le nom
%  du parleur en petites capitales : on garde \textsc{}.
%
%  LES DEUX NOMS PROPRES SE TRADUISENT, ET L'IDO L'A DEJA FAIT :
%  Leblanc est « Blanko », Roux est « Rufo ». Le suedois forme les
%  memes noms sur les memes adjectifs — vit, röd — et ecrit donc herr
%  Vit et herr Rödh, qui sont de vrais noms de famille suedois.
% ===================================================================

%%K t05-apar-1 apar
\VUsaut{6.0mm}
\VUtitre{15.0pt}{60}{TABELL Nr 5}
\VUsaut{1.4mm}
\VUtitre{11.4pt}{0}{Huset och dess byggande.}
\VUsaut{2.2mm}

%%K t05-tit-1 sub
\VUpk{10.2pt}{40}{Ett gott möte.}
\VUsaut{3.0mm}

%%K t05-01-1 p
\textsc{Johan}. --- Farbror, jag vill mycket gärna veta hur man bygger ett
\VUgras{hus}.

%%K t05-01-2 p
\textsc{Farbrodern} (80). --- Det är mycket lätt, min vän, kom med mig, jag ska visa dig det
som herr Bernhard (79) låter bygga och som jag ska bo i.

%%K t05-01-3 p
\textsc{Johan}. --- Vem har då ritat \VUgras{planen} \textsuperscript{(76)} till det här
huset?

%%K t05-01-4 p
\textsc{Farbrodern}. --- \VUgras{Arkitekten} \textsuperscript{(75)}; han lade fram en mycket
tydlig ritning som godkändes, och \VUgras{byggmästaren} \textsuperscript{(77)} började
arbetena.

%%K t05-01-5 p
\textsc{Johan}. --- Vem är då byggmästaren?

%%K t05-01-6 p
\textsc{Farbrodern}. --- Han är herr Vit, far till din vän Jakob. Han leder arbetarna
tillsammans med herr Rödh, arkitekten.

%%K t05-01-7 p
Se! just i rättan tid kommer herr Vit med sin \VUgras{tumstock} \textsuperscript{(78)}, och
herr Rödh med sin plan.

%%K t05-01-8 p
God dag, mina herrar. Hur står det till?

%%K t05-01-9 p
\textsc{Herr Rödh}. --- Mycket bra, tack. God dag, Johan, vad det barnet har vuxit!

%%K t05-01-10 p
\textsc{Farbrodern}. --- Ja, verkligen, han är en liten man. Han är mycket allvarlig, man
måste upplysa honom om allt han ser. I dag skulle han vilja veta hur man bygger ett hus.

%%K t05-tit-2 sub
\VUpk{10.2pt}{40}{Arbetarna.}
\VUsaut{3.0mm}

%%K t05-01-11 p
\textsc{Herr Vit}. --- Nåväl, kom med mig, lille vän, jag ska genast förklara det för dig,
eftersom jag tycker mycket om barn som vill lära sig.

%%K t05-01-12 p
Du ser den arbetaren som håller en \VUgras{spade} \textsuperscript{(82)} i handen: han är en
\VUgras{grovarbetare} \textsuperscript{(81)}. Han gräver ett
\VUgras{dike} \textsuperscript{(46)} för att lägga ner vattenledningen. Han har
också grävt ut marken på den plats där huset står, för att murarn skulle resa de fyra
\VUgras{murarna}. Den del av de murarna som är i marken kallas
\VUgras{grunden} \textsuperscript{(2)}. Utrymmet som ligger mellan grunderna
bildar \VUgras{källaren} \textsuperscript{(3)}. Den källaren har ett litet fönster som kallas
\VUgras{källarfönster} \textsuperscript{(5)}, en
\VUgras{dörr} \textsuperscript{(4)} och en trappa.

%%K t05-01-13 p
\VUgras{Murarn} \textsuperscript{(108)} fortsatte att bygga murarna och lämnade
öppningar för \VUgras{dörrarna} \textsuperscript{(8)} och \VUgras{fönstren}
\textsuperscript{(17)}.

%%K t05-01-14 p
\textsc{Johan}. --- Men den murarn ser jag inte!

%%K t05-01-15 p
Herr \textsc{Vit}. --- Nu har han slutat arbeta i huset. Han står vid
\VUgras{stallet} \textsuperscript{(54)}, på sin \VUgras{ställning}
\textsuperscript{(110)}, han bygger muren till \VUgras{tvättstugan} \textsuperscript{(56)}.

%%K t05-01-16 p
Han har en murslev, ett tråg och ett \VUgras{sänklod} \textsuperscript{(109)}. Men de namnen
kommer du inte att minnas.

%%K t05-01-17 p
\textsc{Johan}. --- Jo visst, för jag ska genast be min farbror om en liten murslev och ett
tråg, och jag ska själv göra ett sänklod med ett snöre.

%%K t05-tit-3 sub
\VUsaut{3.0mm}
\VUpk{10.2pt}{40}{Materialet.}
\VUsaut{3.0mm}

%%K t05-01-18 p
\textsc{Farbrodern}. --- Ha! mycket bra. Men säg mig, Johan, vad behöver man för att bygga ett
hus?

%%K t05-01-19 p
\textsc{Johan}. --- Man behöver \VUgras{huggen sten} \textsuperscript{(59)} och
\VUgras{murbruk} \textsuperscript{(65)}.

%%K t05-01-20 p
\textsc{Farbrodern}. --- Just det, se den mannen, med en stor hatt, på sin
\VUgras{kärra} \textsuperscript{(136)}. Han är en \VUgras{forkarl}
\textsuperscript{(135)}. Varje dag för han hit \VUgras{stenblock} \textsuperscript{(60)}. Han
lastar av sin vagn på gården, och
\VUgras{stenhuggarna} \textsuperscript{(83)} hugger blocken och sågar dem med
sin \VUgras{stensåg} \textsuperscript{(85)}. En \VUgras{handtlangare} \textsuperscript{(146)}
bär dem till murarn som lägger dem på plats.

%%K t05-01-21 p
Under tiden gör \VUgras{brukblandaren} \textsuperscript{(87)} i ordning murbruket. Han behöver
\VUgras{sand} \textsuperscript{(62)}, \VUgras{kalk} \textsuperscript{(63)} och \VUgras{vatten}
\textsuperscript{(64)}. Kalken står bredvid honom i en \VUgras{säck} \textsuperscript{(71)},
en
\VUgras{murarhantlangare} \textsuperscript{(111)} bär sanden till honom på en
\VUgras{skottkärra} \textsuperscript{(112)} och \VUgras{lärlingen}
\textsuperscript{(113)} står vid pumpen (58). Han fyller en \VUgras{hink}
\textsuperscript{(114)}. Murbruket är snart färdigt. \VUgras{Brukbäraren}
\textsuperscript{(107)} tar det och bär det upp, längs stegen, på ställningen, där en murare
arbetar som gör färdigt den sidan av huset.

%%K t05-01-22 p
Och nu, vad kallas den arbetare som gör \VUgras{taket} \textsuperscript{(31)}?

%%K t05-01-23 p
\textsc{Johan}. --- Han är en \VUgras{takläggare} \textsuperscript{(118)}.

%%K t05-tit-4 sub
\VUsaut{3.0mm}
\VUpk{10.2pt}{40}{Husets tak.}
\VUsaut{3.0mm}

%%K t05-01-24 p
Herr \textsc{Vit}. --- Ja, men först kommer \VUgras{timmermannen} \textsuperscript{(115)}.

%%K t05-01-25 p
Timmermannen gör \VUgras{takstolen} \textsuperscript{(35)}. Han fogar samman
\VUgras{bjälkarna} \textsuperscript{(37)} med \VUgras{träpluggar}
\textsuperscript{(117)}. De arbetarna, där uppe, med sina vida sammetsbyxor, är timmermän. Den
ene håller en liten bjälke, och den andre sågar av en plugg med sin \VUgras{yxa}
\textsuperscript{(116)}. Den som står vid \VUgras{skorstenen} \textsuperscript{(41)} är
takläggaren. Han spikar läkt på (*) bjälkarna, och
\VUgras{skiffer} \textsuperscript{(66)} på läkten. Ibland lägger han taktegel.

%%K t05-noto-1 noto
\VUnotes{143.2mm}{%
(*) \VUgras{Balk-o} = T. \textit{Balken}, E. \textit{joist}, F. \textit{solive}, I.
\textit{travicello}, R. \textit{balka}, S. Port. \textit{viga}. Teknisk rot som hittills inte
antagits, men som bör föreslås för att översätta betydelsen av det F. ordet \textit{solive},
som varken är en bjälke F. \textit{poutre}, eller en liten bjälke. Det är en fyrkantig
trästång som bär upp ett golv och ett tak.%
}

%%K t05-01-26 p
Stallets tak är \VUgras{tegeltäckt} \textsuperscript{(55)}, husets är
\VUgras{skiffertäckt} \textsuperscript{(32)} och verandans är av \VUgras{zink}
\textsuperscript{(24)}.

%%K t05-01-27 p
\textsc{Johan}. --- Gör takläggaren också zinktaken?

%%K t05-01-28 p
Herr \textsc{Vit}. --- Nej, men \VUgras{zinkslagaren} \textsuperscript{(121)} eller
\VUgras{blyslagaren}, den som sätter en \VUgras{takglugg} \textsuperscript{(38)} på husets
tak. Han sätter också \VUgras{stuprören} \textsuperscript{(43)} och \VUgras{takrännorna}
\textsuperscript{(42)}. Han löder zinken och plåten. Han har fäst \VUgras{åskledaren}
\textsuperscript{(40)} och \VUgras{vindflöjeln} \textsuperscript{(39)} på
\VUgras{gaveln} \textsuperscript{(36)}.

%%K t05-01-29 p
\textsc{Johan}. --- Är huset då färdigbyggt?

%%K t05-tit-5 sub
\VUsaut{3.0mm}
\VUpk{10.2pt}{40}{Verktygen.}
\VUsaut{3.0mm}

%%K t05-01-30 p
Herr \textsc{Vit}. --- Inte helt. \VUgras{Gipsaren} \textsuperscript{(99)} bygger med
\VUgras{tegelstenar} \textsuperscript{(61)} de mellanväggar som delar det i olika rum. Med
\VUgras{gips} \textsuperscript{(70)} rappar han
\VUgras{taken} \textsuperscript{(26)} och murarna. Nu arbetar han på
\VUgras{verandans} \textsuperscript{(33)} tak. Han har, liksom murarn, en
\VUgras{murslev} \textsuperscript{(100)} och ett \VUgras{tråg}
\textsuperscript{(101)}.

%%K t05-01-31 p
Sedan gör snickaren dörrarna, fönstren, \VUgras{parkettgolven} \textsuperscript{(16)} och
\VUgras{fönsterluckorna} \textsuperscript{(19)}.

%%K t05-01-32 p
Nu arbetar han på gården vid sin \VUgras{hyvelbänk} \textsuperscript{(90)}. Han har en
\VUgras{såg} \textsuperscript{(94)} för att såga träet (69), en
\VUgras{hyvel} \textsuperscript{(91)}, ett \VUgras{skruvstäd}
\textsuperscript{(92)}, ett \VUgras{stämjärn} \textsuperscript{(94 bis)}, en
\VUgras{passare} \textsuperscript{(95)}, och en \VUgras{klubba} av trä
\textsuperscript{(95 bis)}.

%%K t05-01-33 p
\textsc{Johan}. --- Ha! men det är roligare att snickra än att mura. Farbror, du ska köpa en
hyvelbänk, en såg och en hyvel åt mig, för jag ska snart göra en koja åt Medor.

%%K t05-01-34 p
\textsc{Farbrodern}. --- Ja, min gosse.

%%K t05-01-35 p
\textsc{Johan}. --- Vad glad jag är! Och vem är den som står vid ytterdörren?

%%K t05-tit-6 sub
\VUsaut{3.0mm}
\VUpk{10.2pt}{40}{Målningen.}
\VUsaut{3.0mm}

%%K t05-01-36 p
Herr \textsc{Vit}. --- Han är en \VUgras{målare} \textsuperscript{(102)}. Han står på sin
stege med en burk \VUgras{färg} \textsuperscript{(103)} och sin
\VUgras{pensel} \textsuperscript{(104)}. Han målar träet i ytterdörren för att
bevara det längre.

%%K t05-01-37 p
Han klistrar också tapeterna i våningarna.

%%K t05-01-38 p
\VUgras{Tapetseraren} \textsuperscript{(107)} som står på en \VUgras{stege}
\textsuperscript{(106)}, på \VUgras{balkongen} \textsuperscript{(28)}, sätter upp en
\VUgras{markis} \textsuperscript{(29)}.

%%K t05-01-39 p
Arbetaren som står vid det stora fönstret är \VUgras{glasmästaren} \textsuperscript{(96)}. Han
fäster \VUgras{fönsterrutorna} \textsuperscript{(18)} med \VUgras{kitt}
\textsuperscript{(73)}. Den andre arbetaren, som knäböjer vid källardörren är en
\VUgras{låssmed} \textsuperscript{(97)}. Han sätter i ett \VUgras{lås} \textsuperscript{(6)}.
Hans \VUgras{fil} \textsuperscript{(98)} ligger bredvid honom.

%%K t05-01-40 p
\textsc{Johan}. --- Är den arbetare som slår med sin \VUgras{hammare} \textsuperscript{(84)}
på ett järnstycke också en låssmed?

%%K t05-01-41 p
Herr \textsc{Vit}. --- Riktigare sagt, han är en smed (125). Han smider
\VUgras{järnsmiden} \textsuperscript{(67)} åt \VUgras{elektrikern}
\textsuperscript{(124)} som står på \VUgras{vagnslidrets} \textsuperscript{(51)} tak. Han har
en \VUgras{ässja} \textsuperscript{(126)}, ett \VUgras{städ} \textsuperscript{(127)} och en
\VUgras{tång} \textsuperscript{(128)}.

%%K t05-01-42 p
Elektrikern fäster telefonens \VUgras{koppar} \textsuperscript{(44)}
\VUgras{tråd}.

%%K t05-tit-7 sub
\VUpk{10.2pt}{40}{Trädgårdsskötseln.}
\VUsaut{3.0mm}

%%K t05-01-43 p
\textsc{Farbrodern}. --- Längst in på \VUgras{gården} \textsuperscript{(45)}, vid
\VUgras{gallret} \textsuperscript{(47)} och \VUgras{porten} \textsuperscript{(48)} ser du
\VUgras{trädgårdsmästaren} \textsuperscript{(129)}. Han gör i ordning den
\VUgras{lilla trädgården} \textsuperscript{(49)}. Han har grävt om jorden med sin
\VUgras{spade} \textsuperscript{(131)}, han sår frön och krattar med \VUgras{krattan}
\textsuperscript{(130)}. Sedan ska han med sin
\VUgras{vattenkanna} \textsuperscript{(132)} vattna \VUgras{rabatterna}
\textsuperscript{(150)} och plantorna ska växa.

%%K t05-01-44 p
\VUgras{Stalldrängen} \textsuperscript{(133)} som just har skött hästen och gett
den foder gör \VUgras{vagnen} \textsuperscript{(52)} ren med en svamp, en
\VUgras{handpump} \textsuperscript{(134)} och en hink vatten. Sedan ska han
köra in den i vagnslidret igen och stänga dörren med den tjocka \VUgras{regeln}
\textsuperscript{(53)}.

%%K t05-01-45 p
Se så, du är nöjd, tror jag, du har fått tillräckliga förklaringar.

%%K t05-01-46 p
\textsc{Johan}. --- Ja, farbror, men jag skulle gärna se husets insida.

%%K t05-01-47 p
\textsc{Farbrodern}. --- Det ska ske i morgon. Herr Rödh ska förklara planen för dig och nämna
namnet på varje rum.

%%K t05-tit-8 sub
\VUsaut{3.0mm}
\VUpk{10.2pt}{40}{Husets inre indelning.}
\VUsaut{2.4mm}

%%K t05-apar-2 apar
{\centering\textit{(Se planen.)}\par}
\VUsaut{2.4mm}

%%K t05-01-48 p
\textsc{Arkitekten}. --- Kom, Johan, jag ska genast låta dig utforska husets olika delar.

%%K t05-01-49 p
Låt oss gå in på \VUgras{bottenvåningen} \textsuperscript{(7)}. Här är knappen till
\VUgras{ringklockan} \textsuperscript{(11)}. Vi går upp för de tre
\VUgras{trappstegen} \textsuperscript{(14)} på \VUgras{förstukvisten}
\textsuperscript{(9)}, vi går över \VUgras{tröskeln} \textsuperscript{(10)} till
\VUgras{ytterdörren} \textsuperscript{(8)} och här står vi vid foten av
\VUgras{trappan} \textsuperscript{(13)}.

%%K t05-01-50 p
\textsc{Johan}. --- Det är korridoren.

%%K t05-01-51 p
\textsc{Arkitekten}. --- Nej, det är \VUgras{vestibulen} \textsuperscript{(12)}.
\VUgras{Korridoren} \textsuperscript{(a)} är längst bort. Till höger, här är
\VUgras{salongen} \textsuperscript{(b)} och \VUgras{matsalen}
\textsuperscript{(c)}, mitt emot matsalen är \VUgras{köket} \textsuperscript{(d)} och
\VUgras{serveringsrummet} \textsuperscript{(e)}, sedan längst fram \VUgras{arbetsrummet}
\textsuperscript{(f)}. \VUgras{Verandan} \textsuperscript{(23)} med sin \VUgras{glasdörr}
\textsuperscript{(25)} ligger bredvid salongen.

%%K t05-01-52 p
Vi ska genast gå upp för trappan. Håll dig i \VUgras{ledstången} \textsuperscript{(15)} och
räkna trappstegen. De är fjorton. Här är
\VUgras{trappavsatsen} \textsuperscript{(m)}, vi är på första \VUgras{våningen}
\textsuperscript{(27)}.

%%K t05-tit-9 sub
\VUsaut{3.0mm}
\VUpk{10.2pt}{40}{Första våningen.}
\VUsaut{3.0mm}

%%K t05-01-53 p
Mitt emot trappan är \VUgras{avträdet} \textsuperscript{(20)} och
\VUgras{toalettrummet} \textsuperscript{(j)}. Till höger ett vackert
\VUgras{sovrum} \textsuperscript{(g)} med två fönster mot husets
\VUgras{fasad} \textsuperscript{(1)}. Till vänster är \VUgras{badrummet}
\textsuperscript{(1)} och \VUgras{gästrummet} \textsuperscript{(i)} med en stor
\VUgras{alkov} \textsuperscript{(h)}: Bredvid sovrummet är
\VUgras{barnkammaren} \textsuperscript{(k)}. Den ligger vid en rymlig
\VUgras{terrass} \textsuperscript{(21)}. Under den terrassen är ett annat avträde
(20) och, på muren, här är \VUgras{klockan} \textsuperscript{(22)} som ringer till middagen.

%%K t05-01-54 p
\textsc{Johan}. --- Låt oss gå upp på \VUgras{andra våningen}.

%%K t05-01-55 p
\textsc{Arkitekten}. --- Det finns ingen andra våning. Under taket finns bara
\VUgras{vindsrum} \textsuperscript{(33)} och vinden (34). Vi har slutat
utforskningen, vi kan gå ner igen.

%%K t05-01-56 p
\textsc{Johan}. --- Tack, min herre, det är mycket intressant. Jag ska genast berätta för min
far allt vad jag har sett.
\VUsaut{4.0mm}
\VUfilet{22mm}
